% =============================================================================
% Section 07 — Non-Equivalence of Unity and Singularity
% Book I — Ontology of Viability
% Chapter 03 — Unity and Plurality
%
% Purpose:
%   To establish the ontological distinction between unity
%   and numerical singularity, and to examine why neither
%   concept can be reduced to the other.
%
% Conceptual Policy:
%   1. Unity is not equivalent to numerical singularity.
%   2. Numerical oneness does not establish unity of identity.
%   3. Unity does not necessarily imply the existence
%      of exactly one identity.
%   4. Plurality does not necessarily negate unity.
%   5. Unity must not be reduced to indivisibility,
%      homogeneity, or numerical oneness.
%   6. No universal conclusion concerning the number
%      of identities is adopted.
%
% =============================================================================

\section{Non-Equivalence of Unity and Singularity}
\label{sec:non-equivalence-of-unity-and-singularity}

% -----------------------------------------------------------------------------
\subsection{The Problem of Equating Unity with Singularity}

Unity and singularity are frequently treated as
equivalent determinations.

When something is described as one, it may be assumed
that it is singular in every relevant ontological
respect. Conversely, when something is described
as multiple, it may be assumed that unity is absent.

Such assumptions introduce a conceptual reduction
that must be examined.

Unity concerns the ontological determination through
which something is considered unified in a specified
respect.

Singularity, in its numerical sense, concerns
the condition of being one rather than more than
one within a particular domain of enumeration.

These concepts may coincide in particular cases,
but their possible coincidence does not establish
their equivalence.

The question is therefore not whether unity and
singularity can occur together.

It is whether one can be defined, established,
or inferred solely through the other.

The present section argues that such a reduction
is not ontologically justified.

% -----------------------------------------------------------------------------
\subsection{Unity as Distinct from Numerical Oneness}

Unity cannot be defined simply as numerical oneness.

Numerical oneness indicates that something is
counted as one within a specified domain and
under a particular criterion of enumeration.

Unity, however, concerns the ontological respect
in which something is considered unified.

The difference is important because numerical
counting presupposes that the subject of counting
and the criterion of individuation have already
been sufficiently determined.

It does not independently establish the ontological
unity of what is counted.

For example, an existent may be counted as one
object within a particular classification while
the unity of its identity remains an independent
ontological question.

Likewise, a system may be treated as one unit
for a particular practical or analytical purpose
without this designation establishing that
it possesses exactly one identity.

The numerical designation of one therefore
cannot replace an account of unity.

The following distinction must be maintained:

\[
\text{Numerical Oneness}
\neq
\text{Ontological Unity}
\]

This distinction does not deny that numerical
oneness may be relevant to the examination
of unity.

It denies that numerical oneness is sufficient
to establish it.

% -----------------------------------------------------------------------------
\subsection{Singularity as a Determination of Number}

Singularity concerns the numerical determination
of a subject as one within a specified domain
of counting.

Such a determination depends on what is being
counted and on the criterion through which
the counted subject is distinguished.

An existent may be considered singular as
a material object, a system, an organism,
or a bounded configuration.

However, these numerical descriptions do not
necessarily settle the question of identity.

The same existent may be singular under one
criterion of enumeration while involving
a plurality of components, relations,
or manifestations.

Likewise, several entities may be counted
as distinct units without this counting
establishing the ontological nature of their
identities.

Singularity is therefore meaningful only
in relation to a specified subject and
a specified numerical domain.

It cannot be treated as an unrestricted
ontological determination that automatically
applies to every aspect of an existent.

% -----------------------------------------------------------------------------
\subsection{The Difference Between Unity and Indivisibility}

Unity must also be distinguished from
indivisibility.

Indivisibility concerns whether something
can be divided into parts or whether such
division is possible under a specified
criterion.

Unity concerns the ontological respect
in which something is unified.

An existent may be divisible into components
while remaining unified in another respect.

Conversely, something may be treated as
indivisible within a particular analysis
without this establishing the unity of
its identity.

The possibility of division does not
automatically negate unity.

Nor does the absence of an identified
division establish unity.

The distinction can be expressed as follows:

\[
\text{Divisibility}
\not\Rightarrow
\text{Absence of Unity}
\]

and

\[
\text{Indivisibility}
\not\Rightarrow
\text{Unity of Identity}
\]

These expressions prevent the reduction
of unity to the absence of internal parts.

A theory of unity must therefore explain
the relevant ontological relation rather
than rely exclusively on criteria of
division or indivisibility.

% -----------------------------------------------------------------------------
\subsection{The Difference Between Unity and Homogeneity}

Unity is not equivalent to homogeneity.

Homogeneity concerns similarity or sameness
among elements or determinations in a
specified respect.

Unity concerns the ontological unity
attributable to a subject.

An existent may involve heterogeneous
components or differentiated determinations
while remaining unified in a relevant
ontological respect.

Likewise, several homogeneous elements
may remain distinct and may not constitute
one identity.

Similarity does not establish unity,
and difference does not necessarily negate it.

The reduction of unity to homogeneity
would impose the assumption that internal
difference is incompatible with unity.

Such an assumption is not justified.

Unity must therefore remain conceptually
distinct from the degree of similarity
among the elements or determinations
of a subject.

% -----------------------------------------------------------------------------
\subsection{The Possibility of Unity Within Plurality}

Unity and plurality are not necessarily
mutually exclusive.

A subject may be unified in one respect
while involving plurality in another.

An existent may exhibit unity of identity
while containing multiple components,
relations, or determinations.

A composition may be unified as a whole
while its constituents remain plural.

A source may be one while being related
to multiple manifestations.

These possibilities indicate that the
presence of plurality does not automatically
negate every form of unity.

However, the recognition of unity within
plurality must not be used to assume
that every plurality is necessarily
contained within one identity.

The subject and respect of each claim
must remain explicit.

The distinction can be represented
schematically:

\[
\text{Unity in one respect}
\quad+\quad
\text{Plurality in another respect}
\]

This expression indicates a possible
ontological configuration rather than
a universal rule.

It requires that the relevant respects
be distinguished before the relation
between unity and plurality is determined.

% -----------------------------------------------------------------------------
\subsection{The Possibility of Plurality Without Singular Identity}

The non-equivalence of unity and singularity
also requires that plurality not be treated
as a simple negation of every possible
form of unity.

A configuration may involve multiple
identities while exhibiting unity in
another ontological respect.

For example, several identities may participate
in one aggregate or may be related within
a broader configuration.

The unity of that aggregate or configuration
does not necessarily eliminate the distinction
among the identities involved.

Likewise, the existence of multiple identities
does not automatically establish that
they form one higher identity.

The relation between the plurality of
identities and the unity of a broader
configuration must therefore be examined
independently.

This possibility prevents the assumption
that unity can exist only where numerical
singularity is established.

% -----------------------------------------------------------------------------
\subsection{Unity and the Number of Identities}

The distinction between unity and singularity
has a direct consequence for the ontology
of identity.

The unity attributed to an existent cannot
be used, by itself, to establish that
the existent possesses exactly one identity.

Such an inference would confuse the unity
of a subject in a specified respect with
the numerical determination of its identities.

Conversely, the recognition of numerical
multiplicity within an existent does not
automatically establish that multiple
identities are involved.

The number of identities requires an
independent ontological examination.

The following implications must therefore
be rejected as automatic conclusions:

\[
\text{Unity}
\not\Rightarrow
\text{Exactly One Identity}
\]

\[
\text{Numerical Singularity}
\not\Rightarrow
\text{Unity of Identity}
\]

\[
\text{Plurality of Elements}
\not\Rightarrow
\text{Plurality of Identities}
\]

These restrictions preserve the distinction
between the ontological unity of a subject,
its numerical description, and the number
of identities involved.

They also prevent the premature adoption
of a universal doctrine of either one
identity or multiple identities.

% -----------------------------------------------------------------------------
\subsection{Unity and the Problem of the Subject}

Every attribution of unity must specify
its subject.

The statement that something is unified
remains incomplete until it is established
what is being considered unified.

The subject may be:

\begin{itemize}
    \item An existent.
    \item An identity.
    \item A source.
    \item An aggregate.
    \item A composition.
    \item A configuration of relations.
    \item A plurality of manifestations.
\end{itemize}

These subjects may be related, but their
unity cannot be assumed to be interchangeable.

The unity of an aggregate does not automatically
establish the unity of its components.

The unity of a source does not automatically
establish the unity of its manifestations.

The unity of an existent does not automatically
settle the number of identities involved.

The specification of the subject is therefore
a necessary condition for meaningful claims
about unity.

The same requirement applies to numerical
singularity.

A numerical claim is incomplete unless
the domain of enumeration and the criterion
of counting are specified.

% -----------------------------------------------------------------------------
\subsection{The Limits of Numerical Classification}

Numerical classification may be useful
for distinguishing one subject from several
subjects within a specified domain.

However, numerical classification cannot
provide a complete ontological account
of unity or identity.

Counting presupposes a criterion through
which the relevant units are identified.

That criterion may be appropriate for
a particular domain while remaining
insufficient for determining identity.

A collection may be counted as one aggregate,
a system as one functional unit, or an
existent as one material object.

None of these numerical classifications
automatically determines whether one
or multiple identities are involved.

The ontological significance of numerical
classification must therefore be established
rather than assumed.

A numerical description may contribute
to an ontological inquiry, but it cannot
replace the distinctions required by
that inquiry.

% -----------------------------------------------------------------------------
\subsection{Unity Without an Assumption of Universal Singularity}

Rejecting the equivalence between unity
and singularity does not mean rejecting
the possibility of singular identity.

An existent may possess one identity
in a particular ontological configuration.

Such a conclusion, however, requires
an argument concerning identity itself.

It cannot be established merely through
the numerical designation of one existent
or through the recognition of unity in
its composition, source, or manifestation.

Likewise, the possibility of multiple
identities must remain open where the
relevant ontological conditions support
such an examination.

The present section therefore does not
replace a universal doctrine of singularity
with a universal doctrine of plurality.

It rejects the assumption that the question
can be settled through the concept of unity
alone.

The number and relation of identities
must be examined independently.

% -----------------------------------------------------------------------------
\subsection{Conclusion: Unity and Singularity as Distinct Concepts}

Unity and singularity are distinct ontological
and numerical determinations.

Unity concerns the ontological respect
in which a subject is considered unified.

Singularity concerns its numerical
determination as one within a specified
domain of enumeration.

Neither concept can be reduced to the other.

Unity does not automatically establish
exactly one identity.

Numerical singularity does not automatically
establish unity of identity.

Divisibility, indivisibility, homogeneity,
and numerical classification likewise
cannot serve as exhaustive definitions
of unity.

The distinction also makes it possible
to recognize unity and plurality within
the same broader ontological configuration
without treating them as contradictory
in every respect.

The next section examines the limits
of unity and plurality.

It considers the conditions under which
claims of unity or plurality become
ontologically insufficient, and why
neither determination can be extended
beyond the subject and respect for which
it has been established.

% =============================================================================
% End of Section 07
% =============================================================================